\documentclass[11pt,a4paper]{article}

\usepackage[T1]{fontenc}
\usepackage[utf8]{inputenc}
\usepackage[french]{babel}
\usepackage{lmodern}
\usepackage{microtype}
\usepackage{geometry}
\geometry{margin=2.5cm}

\usepackage{amsmath,amssymb,amsthm,mathtools}
\usepackage{bbm}
\usepackage{enumitem}
\usepackage{hyperref}
\hypersetup{colorlinks=true,linkcolor=blue,urlcolor=blue,citecolor=blue}

\newcommand{\E}{\mathbb{E}}
\newcommand{\Pp}{\mathbb{P}}
\newcommand{\R}{\mathbb{R}}
\newcommand{\N}{\mathbb{N}}
\newcommand{\1}{\mathbbm{1}}
\newcommand{\cF}{\mathcal{F}}
\newcommand{\cA}{\mathcal{A}}
\newcommand{\cB}{\mathcal{B}}
\newcommand{\Lip}{\mathrm{Lip}}
\newcommand{\Var}{\mathrm{Var}}
\newcommand{\Cov}{\mathrm{Cov}}

\theoremstyle{plain}
\newtheorem{theorem}{Théorème}
\newtheorem{proposition}{Proposition}
\newtheorem{lemma}{Lemme}
\newtheorem{corollary}{Corollaire}

\theoremstyle{definition}
\newtheorem{definition}{Définition}
\newtheorem{remark}{Remarque}

\title{Probabilités numériques / Méthodes de Monte Carlo -- Corrigés (Sujets 07 à 10) \\
\large et fiche de révision (définitions, théorèmes, méthodes)}
\author{}
\date{Décembre 2025}

\begin{document}
\maketitle
\tableofcontents
\bigskip

\section*{Avant-propos}
Ce document regroupe une correction rigoureuse et rédigée de quatre sujets:
\begin{itemize}[leftmargin=*]
  \item \textbf{Sujet 07} (3 janvier 2011): sensibilités (différences finies et accélération), digitales, réduction de variance (importance sampling gaussien).
  \item \textbf{Sujet 08} (2 avril 2012): équidistribution et QMC (van der Corput), co-monotonie et vega, convergence p.s. d'Euler, volatilité locale et propagation de convexité.
  \item \textbf{Sujet 09} (10 janvier 2013): schéma de Milstein, positivité, modèles CIR/OU, Heston, ponts de diffusion et barrières.
  \item \textbf{Sujet 10} (5 janvier 2015): quasi-Monte Carlo, randomisation, et estimation de la sensibilité d'une digitale (``kappa'').
\end{itemize}
Les notations suivent celles usuelles du cours: $W$ désigne un mouvement brownien standard; pour un pas de discrétisation $\Delta=T/n$ on note $t_i=i\Delta$ et $\Delta W_i = W_{t_{i+1}}-W_{t_i}\sim\mathcal{N}(0,\Delta)$.

\part{Sujet 07 (3 janvier 2011)}

\section{Problème 1 -- Sensibilités I: accélérer les différences finies}
\subsection{Cadre et notations}
Soit $F:\R\times\R^q\to\R$ borélienne et $Z$ un vecteur aléatoire. Fixons $x\in\R$ et $\varepsilon_0>0$ tels que $F(x,Z)\in L^2$ et
\begin{equation}\label{eq:L2-Lip-F}
  \forall x'\in (x-\varepsilon_0,x+\varepsilon_0),\qquad \|F(x',Z)-F(x,Z)\|_{L^2}\le C_{F,Z}|x'-x|.
\end{equation}
On pose $f(\xi)=\E[F(\xi,Z)]$ et on suppose $f\in C^3$ sur $I_x=(x-\varepsilon_0,x+\varepsilon_0)$, avec $f^{(3)}$ Lipschitz de constante $[f^{(3)}]_{\Lip}$ sur $I_x$.

\subsection{1 -- Développement de la différence finie symétrique}
\begin{proposition}\label{prop:df-sym-ordre2}
Pour tout $\varepsilon\in(0,\varepsilon_0)$,
\[
\left|\frac{f(x+\varepsilon)-f(x-\varepsilon)}{2\varepsilon}-f'(x)-\frac{\varepsilon^2}{6}f^{(3)}(x)\right|
\le \frac{\varepsilon^3}{6}[f^{(3)}]_{\Lip}.
\]
\end{proposition}
\begin{proof}
On applique Taylor à l'ordre 3 avec reste intégral (ou forme de Lagrange contrôlée par la Lipschitzité de $f^{(3)}$):
\[
f(x\pm\varepsilon)
=f(x)\pm \varepsilon f'(x)+\frac{\varepsilon^2}{2}f''(x)\pm \frac{\varepsilon^3}{6}f^{(3)}(x)+R_\pm,
\]
avec $|R_\pm|\le \frac{\varepsilon^4}{24}[f^{(3)}]_{\Lip}$.
En soustrayant $f(x-\varepsilon)$ de $f(x+\varepsilon)$ et en divisant par $2\varepsilon$, les termes pairs s'annulent, on obtient
\[
\frac{f(x+\varepsilon)-f(x-\varepsilon)}{2\varepsilon}
=f'(x)+\frac{\varepsilon^2}{6}f^{(3)}(x)+\frac{R_+-R_-}{2\varepsilon},
\]
et $|(R_+-R_-)/(2\varepsilon)|\le (\varepsilon^3/6)[f^{(3)}]_{\Lip}$.
\end{proof}

\subsection{2 -- Estimateur MC de la différence finie}
Pour $\varepsilon\in(0,\varepsilon_0)$, on définit
\[
F_\varepsilon(x,z)=\frac{F(x+\varepsilon,z)-F(x-\varepsilon,z)}{2\varepsilon},
\qquad
\widehat f'_{\varepsilon,M}(x)=\frac1M\sum_{k=1}^M F_\varepsilon(x,Z_k),
\]
où $(Z_k)_{k\ge1}$ sont des copies i.i.d. de $Z$.

\subsubsection{2.a -- Espérance, LLN et CLT}
\begin{proposition}
On a $\E[F_\varepsilon(x,Z)]=(f(x+\varepsilon)-f(x-\varepsilon))/(2\varepsilon)$.
De plus, quand $M\to\infty$,
\[
\widehat f'_{\varepsilon,M}(x)\xrightarrow[]{a.s.}\E[F_\varepsilon(x,Z)],
\qquad
\sqrt{M}\Big(\widehat f'_{\varepsilon,M}(x)-\E[F_\varepsilon(x,Z)]\Big)\xRightarrow[]{}
\mathcal N(0,\sigma_{f,\varepsilon}^2),
\]
où $\sigma_{f,\varepsilon}^2=\Var(F_\varepsilon(x,Z))$.
Enfin, sous \eqref{eq:L2-Lip-F}, on a la borne $\sigma_{f,\varepsilon}^2\le C_{F,Z}^2$.
\end{proposition}
\begin{proof}
L'identité d'espérance est immédiate par linéarité de $\E$.
La convergence p.s. et le TCL suivent du théorème de la loi forte et du TCL i.i.d. dès que $F_\varepsilon(x,Z)\in L^2$, ce qui découle de \eqref{eq:L2-Lip-F}.
Enfin
\[
\|F_\varepsilon(x,Z)\|_{L^2}
=\frac{1}{2\varepsilon}\|F(x+\varepsilon,Z)-F(x-\varepsilon,Z)\|_{L^2}
\le \frac{1}{2\varepsilon}\,C_{F,Z}\,|2\varepsilon|=C_{F,Z},
\]
donc $\Var(F_\varepsilon(x,Z))\le \E[F_\varepsilon(x,Z)^2]\le C_{F,Z}^2$.
\end{proof}

\subsubsection{2.b -- Contrôle de l'erreur quadratique}
\begin{proposition}\label{prop:mse-df-2011}
Pour tout $\varepsilon\in(0,\varepsilon_0)$ et $M\ge 1$,
\[
\E\Big[\big(f'(x)-\widehat f'_{\varepsilon,M}(x)\big)^2\Big]
\le
\left(\frac{\varepsilon^2}{6}|f^{(3)}(x)|+\frac{\varepsilon^3}{6}[f^{(3)}]_{\Lip}\right)^2
\;+\;\frac{\sigma_{f,\varepsilon}^2}{M}.
\]
\end{proposition}
\begin{proof}
On décompose l'erreur en \emph{biais} + \emph{fluctuation}:
\[
f'(x)-\widehat f'_{\varepsilon,M}(x)
\;=\;\Big(f'(x)-\E[F_\varepsilon(x,Z)]\Big)\;+\;\Big(\E[F_\varepsilon(x,Z)]-\widehat f'_{\varepsilon,M}(x)\Big).
\]
Le second terme est centré; en développant le carré on obtient
\[
\E[(f'(x)-\widehat f'_{\varepsilon,M}(x))^2]
=\big(f'(x)-\E[F_\varepsilon(x,Z)]\big)^2+\Var(\widehat f'_{\varepsilon,M}(x)).
\]
Le biais est contrôlé par la proposition \ref{prop:df-sym-ordre2}. Enfin
$\Var(\widehat f'_{\varepsilon,M})=\Var(F_\varepsilon)/M$.
\end{proof}

\subsubsection{2.c -- Choix de $\varepsilon$ et $M$ (biais/variance)}
La borne \ref{prop:mse-df-2011} montre ici un point important:
\begin{itemize}[leftmargin=*]
  \item le \textbf{biais} est d'ordre $\varepsilon^2$ (donc biais$^2$ d'ordre $\varepsilon^4$);
  \item la \textbf{variance MC} est d'ordre $1/M$ et, sous \eqref{eq:L2-Lip-F}, \emph{ne dépend pas} (au pire) de $\varepsilon$.
\end{itemize}
Ainsi, à $M$ fixé, diminuer $\varepsilon$ améliore essentiellement le biais sans pénaliser la variance (contrairement aux digitales).

\subsection{3 -- Accélération (Richardson) de la différence finie}
\subsubsection{3.a -- Combinaison qui annule le terme en $\varepsilon^2$}
\begin{proposition}\label{prop:richardson-2011}
Il existe $a,b\in\R$ (uniques) tels que $a+b=1$ et, pour tout $\varepsilon\in(0,\varepsilon_0)$,
\[
\left|\E\big[aF_\varepsilon(x,Z)+bF_{\varepsilon/2}(x,Z)\big]-f'(x)\right|
\le \frac{\varepsilon^3}{12}[f^{(3)}]_{\Lip}.
\]
On peut prendre $a=-\frac13$ et $b=\frac43$.
\end{proposition}
\begin{proof}
Par la proposition \ref{prop:df-sym-ordre2},
\[
\E[F_\varepsilon(x,Z)]=f'(x)+\frac{\varepsilon^2}{6}f^{(3)}(x)+O(\varepsilon^3),
\qquad
\E[F_{\varepsilon/2}(x,Z)]=f'(x)+\frac{\varepsilon^2}{24}f^{(3)}(x)+O(\varepsilon^3).
\]
Choisir $a,b$ tels que $a+b=1$ et $\frac{a}{6}+\frac{b}{24}=0$ annule le terme en $\varepsilon^2$.
On trouve $b=-4a$ puis $a=-1/3$, $b=4/3$.
Le reste est d'ordre $\varepsilon^3$ et se contrôle explicitement, d'où la borne.
\end{proof}

\subsubsection{3.b -- Estimateur accéléré et contrôle MSE}
On définit l'estimateur accéléré (biaisé) par
\[
\widehat f'^{\,g}_{\varepsilon,M}(x)
=\frac1M\sum_{k=1}^M\left(-\frac13F_\varepsilon(x,Z_k)+\frac43F_{\varepsilon/2}(x,Z_k)\right).
\]
Le biais est désormais d'ordre $\varepsilon^3$ (biais$^2$ d'ordre $\varepsilon^6$).
Pour la variance, l'astuce importante est d'utiliser les \textbf{mêmes} $Z_k$ (common random numbers) pour $F_\varepsilon$ et $F_{\varepsilon/2}$ afin de bénéficier de corrélations favorables.

\subsubsection{3.c -- Astuce de réordonnancement (amélioration de constante)}
En réécrivant la combinaison linéaire en regroupant les évaluations de $F$:
\[
-\frac13F_\varepsilon+\frac43F_{\varepsilon/2}
=\frac{1}{3\varepsilon}\Big(-F(x+\varepsilon,\cdot)+F(x-\varepsilon,\cdot)+4F(x+\varepsilon/2,\cdot)-4F(x-\varepsilon/2,\cdot)\Big),
\]
on peut majorer le $L^2$ de façon plus fine que via l'inégalité triangulaire brute (ce qui améliore la constante de variance asymptotique, conformément à l'énoncé).

\subsubsection{3.d -- Choix $\varepsilon, M$}
Avec l'accélération, le biais$^2$ est $\asymp \varepsilon^6$ tandis que la variance reste $\asymp 1/M$ (au pire).
Pour diviser la MSE par 2, on peut typiquement:
\begin{itemize}[leftmargin=*]
  \item doubler $M$ (variance divisée par 2);
  \item réduire $\varepsilon$ d'un facteur $2^{1/6}$ (biais$^2$ divisé par 2).
\end{itemize}
On constate que la régularité supplémentaire exploitée (Richardson) permet de \emph{réduire beaucoup plus lentement} $\varepsilon$ pour atteindre le même gain de précision.

\section{Problème 2 -- Sensibilités II: options digitales}
\subsection{Cadre}
On considère
\[
dX_t^x=rX_t^x\,dt+\sigma(X_t^x)\,dW_t,\qquad X_0^x=x>0,
\]
où $\sigma:\R\to\R_+^*$ est bornée, dérivable et $\sup_{\xi\in\R}|\xi\sigma'(\xi)|\le C$.

\subsection{1 -- Existence, représentation exponentielle et positivité}
\paragraph{1.a.}
Sur $(0,\infty)$, les coefficients sont localement Lipschitz (sous les hypothèses données) et à croissance linéaire, donc il existe une solution forte unique (théorème standard; la positivité ci-dessous évite toute atteinte de $0$).
\paragraph{1.b--1.c.}
On applique Itô à $\log X_t^x$ (possible tant que $X_t^x>0$) ou bien on reconnaît $X$ comme une exponentielle de Doléans:
\[
X_t^x
=x\exp\Big(\int_0^t\Big(r-\frac12\Big(\frac{\sigma(X_s^x)}{X_s^x}\Big)^2\Big)ds+\int_0^t \frac{\sigma(X_s^x)}{X_s^x}\,dW_s\Big).
\]
Le membre de droite est bien défini (les intégrandes sont contrôlées via la bornitude de $\sigma$ et les estimées de moments), et la formule d'Itô montre qu'il résout l'EDS; en particulier c'est une exponentielle, donc $X_t^x>0$ p.s.
En posant $Y_t^{(x)}=\log X_t^x$, on obtient
\[
dY_t^{(x)}=\Big(r-\frac12\sigma^2(e^{Y_t^{(x)}})\Big)dt+\sigma(e^{Y_t^{(x)}})\,dW_t,\qquad Y_0^{(x)}=\log x.
\]

\subsection{2 -- Dérivabilité de $f(x)=\Pp(X_T^x\ge K)$}
Sous les hypothèses d'ellipticité uniforme et de régularité admises dans l'énoncé, $Y_T^{(x)}$ admet une densité $q(T,x,\cdot)$ $C^\infty$ et avec bornes sur les dérivées.
Comme
\[
f(x)=\Pp(Y_T^{(x)}\ge \log K)=\int_{\log K}^{\infty} q(T,x,y)\,dy,
\]
on peut dériver sous le signe intégral et obtenir
\[
f'(x)=\int_{\log K}^{\infty}\partial_x q(T,x,y)\,dy,
\]
et plus généralement $f\in C^\infty$.

\subsection{2.b--2.c -- Astuce monotonicité $\Rightarrow$ Hölder $1/2$ en $L^2$}
On utilise le rappel: si $x'\ge x$ alors $X_t^{x'}\ge X_t^{x}$ p.s. pour tout $t$.
Ainsi, pour $x<x'$,
\[
\big(\1_{\{X_T^{x}\ge K\}}-\1_{\{X_T^{x'}\ge K\}}\big)^2
=\1_{\{X_T^{x}<K\le X_T^{x'}\}}.
\]
En prenant l'espérance,
\[
\E\Big[\big(\1_{\{X_T^{x}\ge K\}}-\1_{\{X_T^{x'}\ge K\}}\big)^2\Big]
=\Pp(X_T^{x'}\ge K)-\Pp(X_T^{x}\ge K)=f(x')-f(x).
\]
Sur un compact $I\Subset(0,\infty)$, $f'$ est bornée, donc $f$ est Lipschitz sur $I$ et
\[
\big\|\1_{\{X_T^{x}\ge K\}}-\1_{\{X_T^{x'}\ge K\}}\big\|_{L^2}
=\sqrt{f(x')-f(x)}\ \le\ C_I^{1/2}|x'-x|^{1/2}.
\]
Autrement dit, la fonction aléatoire $x\mapsto \1_{\{X_T^{x}\ge K\}}$ est Hölder d'ordre $1/2$ à valeurs dans $L^2$.

\subsection{3 -- Différences finies pour la digitale: explosion de variance}
On pose $F(x,\omega)=\1_{\{X_T^{x}(\omega)\ge K\}}$ et, pour $\varepsilon>0$ petit,
\[
F_\varepsilon(x,\omega)=\frac{F(x+\varepsilon,\omega)-F(x-\varepsilon,\omega)}{2\varepsilon},
\qquad
\widehat f'_{\varepsilon,M}(x)=\frac1M\sum_{k=1}^M F_\varepsilon(x,\omega_k).
\]
Comme $F$ n'est que Hölder $1/2$ en $L^2$, on a typiquement
\[
\|F(x+\varepsilon,\cdot)-F(x-\varepsilon,\cdot)\|_{L^2}\le C\sqrt{\varepsilon},
\]
ce qui implique une borne de variance
\[
\Var(F_\varepsilon(x,\cdot))\ \lesssim\ \frac{1}{\varepsilon}.
\]
Conclusion (astuce à retenir): \textbf{pour les digitales}, diminuer $\varepsilon$ réduit le biais mais \emph{augmente} la variance (phénomène absent au Problème 1).
L'optimisation biais/variance mène alors à des choix de type $\varepsilon\asymp M^{-1/5}$ (ordre classique).
\medskip

\paragraph{4. Amélioration.}
On peut reprendre l'accélération de Richardson du Problème 1 pour réduire l'ordre du biais (passer de $\varepsilon^2$ à $\varepsilon^3$), mais la variance gardant une dépendance en $1/\varepsilon$, l'optimum change et l'amélioration est réelle mais plus délicate: il faut \emph{re-optimiser} simultanément $\varepsilon$ et $M$.

\section{Problème 3 -- Réduction de variance: importance sampling gaussien}
\subsection{Cadre}
Soit $Z\sim\mathcal N(0,I_d)$ et $\phi(Z)=(K-\sum_{k=1}^d \lambda_k e^{\sigma_k Z_k})_+$.
On veut estimer $\E[\phi(Z)]$.

\subsection{1 -- Identité de Cameron--Martin (translation gaussienne)}
\begin{proposition}\label{prop:cameron-martin}
Pour tout $\theta\in\R^d$,
\[
\E[\phi(Z)]=\E\Big[\phi(Z+\theta)\exp\Big(-\langle\theta,Z\rangle-\frac{\|\theta\|^2}{2}\Big)\Big].
\]
\end{proposition}
\begin{proof}
Il s'agit de la formule de changement de mesure/translation pour la loi gaussienne (Cameron--Martin). On peut aussi la vérifier par un calcul direct de densités.
\end{proof}

\subsection{2 -- Estimateur et fonction $g(\theta)$}
L'estimateur MC associé à \ref{prop:cameron-martin} est
\[
\widehat I_M(\theta)=\frac1M\sum_{m=1}^M \phi(Z^{(m)}+\theta)\exp\Big(-\langle\theta,Z^{(m)}\rangle-\frac{\|\theta\|^2}{2}\Big).
\]
La quantité
\[
g(\theta)=\E\Big[\phi(Z+\theta)^2\exp\big(-2\langle\theta,Z\rangle-\|\theta\|^2\big)\Big]
\]
est (à constante près) le \textbf{second moment} de l'estimateur et donc contrôle sa variance; minimiser $g(\theta)$ revient à minimiser la variance.

\subsection{3 -- Convexité stricte et existence d'un minimiseur unique}
On réécrit $g$ sous une forme équivalente (changement de variable $Z\mapsto Z-\theta$):
\[
g(\theta)=\E\Big[\phi(Z)^2\exp\Big(-\langle\theta,Z\rangle+\frac{\|\theta\|^2}{2}\Big)\Big].
\]
Comme $\Pp(\phi(Z)>0)>0$ (en particulier dans le régime $K<\sum_k\lambda_k$ étudié plus loin), l'intégrande est strictement convexe en $\theta$ et $g(\theta)\to\infty$ lorsque $\|\theta\|\to\infty$, ce qui assure l'existence et l'unicité de
\[
\theta^*=\arg\min_{\theta\in\R^d} g(\theta).
\]
De plus, pour $k=1,\dots,d$,
\[
\frac{\partial g}{\partial\theta_k}(\theta)
=\E\Big[\phi(Z)^2\exp\Big(-\langle\theta,Z\rangle+\frac{\|\theta\|^2}{2}\Big)\,(\theta_k-Z_k)\Big].
\]
(Astuce: cette formule permet une optimisation stochastique type Robbins--Monro.)

\subsection{4--6 -- Choix explicite $\bar\theta$ et interprétation}
L'inclusion de niveau issue de la concavité du $\log$ (et de Jensen) permet de montrer qu'un choix de la forme
\[
\theta(\alpha)=\alpha(\lambda_1\sigma_1,\dots,\lambda_d\sigma_d)
\]
donne une direction de descente pour $g$ dès que $\alpha$ est assez grand, et en particulier que le choix explicite $\bar\theta$ de l'énoncé satisfait $g(\bar\theta)<g(0)$.
Conclusion: il est pertinent d'utiliser \textbf{importance sampling gaussien} (shift) pour estimer un put panier, en initialisant le shift à $\bar\theta$ puis en optimisant éventuellement numériquement.

\part{Sujet 08 (2 avril 2012)}

\section{Exercice 1 -- Équidistribution, van der Corput et QMC}
\subsection{1.a -- Définition et caractérisations}
\begin{definition}[Équidistribution sur $[0,1]^d$]
Une suite $(\xi_n)_{n\ge1}\subset[0,1]^d$ est dite équirepartie (équidistribuée) si pour tout pavé $B=\prod_{j=1}^d [0,a_j)$,
\[
\frac1n\sum_{k=1}^n \1_B(\xi_k)\xrightarrow[n\to\infty]{} \lambda_d(B)=\prod_{j=1}^d a_j.
\]
\end{definition}
\begin{remark}[Critère de Weyl]
Équivalent: pour tout $h\in\Z^d\setminus\{0\}$, $\frac1n\sum_{k=1}^n e^{2\pi i\langle h,\xi_k\rangle}\to 0$.
\end{remark}

\subsection{1.b -- Discrépance faible}
On dit que $(\xi_n)$ est \emph{à discrépance faible} si $D_n^*(\xi_1,\dots,\xi_n)\to 0$ (et typiquement avec un taux proche de $(\log n)^d/n$).

\subsection{2 -- Suite de van der Corput en base $p$}
Soit $p\ge2$ et $n=\sum_{k\ge 0} a_k p^k$ sa décomposition $p$-adique.
On définit
\[
\xi_n=\sum_{k\ge0}\frac{a_k}{p^{k+1}}.
\]
\subsubsection{2.a}
Pour $r\in\{0,\dots,p-1\}$, l'écriture $np+r$ décale les chiffres: on obtient
\[
\xi_{np+r}=\frac{\xi_n+r}{p}.
\]
Comme $(\xi_n)$ est équirepartie, pour tout $r$ fixé, la sous-suite $(\xi_{kp+r})_{k\ge1}$ est équirepartie sur $[0,1]$ avec transformation affine, et
\[
\lim_{n\to\infty}\frac1n\sum_{k=1}^n \xi_{kp+r}
=\int_0^1 \frac{u+r}{p}\,du
=\frac{r+1/2}{p}.
\]
\subsubsection{2.b}
De même, pour $r,r'\in\{0,\dots,p-1\}$,
\[
\lim_{n\to\infty}\frac1n\sum_{k=1}^n \xi_{kp+r}\xi_{kp+r'}
=\int_0^1 \frac{u+r}{p}\frac{u+r'}{p}\,du
=\frac{1}{p^2}\Big(\frac13+\frac{r+r'}{2}+rr'\Big).
\]
En soustrayant le produit des limites des moyennes, on obtient
\[
\lim_{n\to\infty}\left(\frac1n\sum_{k=1}^n \xi_{kp+r}\xi_{kp+r'}-\Big(\frac1n\sum_{k=1}^n \xi_{kp+r}\Big)\Big(\frac1n\sum_{k=1}^n \xi_{kp+r'}\Big)\right)
=\frac{1}{12p^2},
\]
indépendant de $r,r'$.

\subsubsection{3 -- Commentaire (astuce/prudence)}
Ce calcul met en évidence une \emph{structure arithmétique forte} et des corrélations non négligeables entre sous-suites.
Moralité QMC à retenir: en dimension $d>1$, il faut utiliser des constructions adaptées (nets/suites Sobol', Halton multi-d) ou des \emph{scramblings}, plutôt que des couplages naïfs de sous-suites très dépendantes.

\subsection{3.a--3.b -- QMC pour $\E[f(Z)]$, $Z\sim\mathcal N(0,\Sigma)$}
\begin{itemize}[leftmargin=*]
  \item Simuler $U_k\in[0,1]^2$ via une suite à faible discrépance; poser $G_k=(\Phi^{-1}(U_k^{(1)}),\Phi^{-1}(U_k^{(2)}))\sim \mathcal N(0,I_2)$ au sens ``transport''.
  \item Pour une covariance $\Sigma$, prendre une factorisation $\Sigma=AA^\top$ (Cholesky) et poser $Z_k=AG_k$.
  \item Approximer $\E[f(Z)]\approx \frac1n\sum_{k=1}^n f(Z_k)$, et randomiser (shift/scramble) pour obtenir des IC.
\end{itemize}

\section{Exercice 2 -- Co-monotonie et application à la vega}
\subsection{2.a--2.c -- Inégalité de Chebyshev (covariance positive)}
\begin{proposition}[Co-monotonie]\label{prop:comono}
Si $f$ et $g$ sont monotones de même sens et $f(Z),g(Z)$ intégrables, alors
\[
\E[f(Z)g(Z)]\ge \E[f(Z)]\,\E[g(Z)].
\]
\end{proposition}
\begin{proof}
Soient $Z'$ une copie indépendante de $Z$. Alors $(f(Z)-f(Z'))(g(Z)-g(Z'))\ge 0$ p.s. par co-monotonie.
En prenant l'espérance et en développant,
\[
0\le \E[(f(Z)-f(Z'))(g(Z)-g(Z'))]=2\Big(\E[f(Z)g(Z)]-\E[f(Z)]\E[g(Z)]\Big).
\]
\end{proof}
\begin{remark}
Pour des fonctions positives non nécessairement dans $L^2$, on conclut par troncature et convergence monotone.
\end{remark}

\subsection{3 -- Vega en Black--Scholes et signe}
Dans Black--Scholes: $X_T=x_0\exp((r-\sigma^2/2)T+\sigma W_T)$, et
\[
\Phi(x_0,\sigma,r,T)=e^{-rT}\E[\varphi(X_T)].
\]
\paragraph{3.a (astuce).}
Pour $\varphi$ convexe $C^1$ à croissance polynomiale, on obtient des inégalités de pentes
\[
\varphi(x)-\varphi(x-1)\le \varphi'(x)\le \varphi(x+1)-\varphi(x),
\]
ce qui assure la croissance polynomiale de $\varphi'$ et justifie la différentiation sous l'espérance.
\paragraph{3.b.}
Comme $\partial_\sigma X_T=X_T(W_T-\sigma T)$, on a
\[
\frac{\partial\Phi}{\partial\sigma}(x_0,\sigma,r,T)=e^{-rT}\E\big[\varphi'(X_T)\,X_T\,(W_T-\sigma T)\big].
\]
\paragraph{3.c.}
En effectuant un changement de mesure de Girsanov qui recentre $W_T-\sigma T$ (ou une réécriture équivalente), on peut ramener la formule à une espérance sous $Z\sim\mathcal N(0,1)$:
\[
\frac{\partial\Phi}{\partial\sigma}(x_0,\sigma,r,T)=x_0\sqrt{T}\,\E\big[\varphi'(x_0e^{\sigma\sqrt{T}Z+(r+\sigma^2/2)T})\,Z\big].
\]
\paragraph{3.d (signe).}
Si $\varphi$ est convexe, alors $\varphi'$ est croissante; la fonction $z\mapsto z$ est croissante.
Par co-monotonie (prop. \ref{prop:comono}) et $\E[Z]=0$, on obtient $\E[\varphi'(\cdot)\,Z]\ge 0$, donc \(\partial_\sigma\Phi\ge 0\).

\section{Exercice 3 -- Convergence p.s. du schéma d'Euler}
\subsection{3.a--3.b -- De $L^p$ à p.s. (Borel--Cantelli)}
Sous hypothèses Lipschitz, on sait (cours) que, pour tout $p\ge 1$,
\[
\E\Big[\sup_{t\le T}|X_t-\bar X_t^n|^p\Big]\le C_p\,n^{-p/2}.
\]
Fixons $\alpha\in(0,1/2)$. Choisissons $p$ tel que $p(1/2-\alpha)>1$. Alors
\[
\sum_{n\ge 1}\Pp\Big(n^\alpha\sup_{t\le T}|X_t-\bar X_t^n|>\varepsilon\Big)
\le \sum_{n\ge 1}\frac{n^{p\alpha}}{\varepsilon^p}\E\Big[\sup_{t\le T}|X_t-\bar X_t^n|^p\Big]
\le C\sum_{n\ge 1} n^{-p(1/2-\alpha)}<\infty.
\]
Par Borel--Cantelli, $n^\alpha\sup_{t\le T}|X_t-\bar X_t^n|\to 0$ p.s.
\subsection{3.c}
Le même type de résultat vaut pour l'Euler constant par morceaux, en combinant l'erreur de discrétisation et l'oscillation intra-pas (modulus de continuité du brownien).

\section{Problème -- Volatilité locale: positivité, convexité, comparaison}
\subsection{1 -- Positivité et SDE du logarithme}
Pour $dX_t^x=X_t^x\sigma(X_t^x)\,dW_t$ (prix actualisé) avec $\sigma$ bornée, Lipschitz sur $\R_+$, on obtient
\[
X_t^x=x\exp\Big(-\frac12\int_0^t \sigma^2(X_s^x)\,ds+\int_0^t \sigma(X_s^x)\,dW_s\Big),
\]
donc $X_t^x>0$ p.s. et $\xi_t=\log X_t^x$ vérifie
\[
d\xi_t=-\frac12\sigma^2(e^{\xi_t})\,dt+\sigma(e^{\xi_t})\,dW_t.
\]

\subsection{2 -- Outils convexes (astuces)}
Pour $g$ convexe Lipschitz, ses dérivées à gauche/droite $g'_-,g'_+$ existent partout, sont croissantes et bornées.
Un outil clé: si $Z$ est centrée sans atome et $g$ convexe Lipschitz croissante, alors
\[
Q(g)(x,u)=\E[g(x+uZ)]
\]
est dérivable en $u>0$ avec
\[
\partial_u Q(g)(x,u)=\E[g'_+(x+uZ)\,Z],
\]
et $Q(g)$ est croissante en $u$, convexe, Lipschitz en $(x,u)$.

\subsection{3--5 -- Schéma d'Euler, opérateur $Q_+$ et propagation}
Pour l'Euler discret, on écrit (avec $Z_k\sim\mathcal N(0,1)$ i.i.d.)
\[
\bar X_{t_k}^n=\bar X_{t_{k-1}}^n\Big(1+\sigma(\bar X_{t_{k-1}}^n)Z_k\sqrt{T/n}\Big),
\]
et $(\bar X_{t_k}^n)_+$ est une chaîne de Markov.
En posant
\[
\bar M_k^n=\E[\varphi((\bar X_T^n)_+)\mid \sigma(Z_1,\dots,Z_k)],
\]
on montre par récurrence qu'il existe des fonctions $\varphi_k$ telles que $\bar M_k^n=\varphi_k((\bar X_{t_k}^n)_+)$ et
\[
\varphi_k(x)=Q_+(\varphi_{k+1})\Big(x,\;x\sigma(x)\sqrt{T/n}\Big),
\qquad \varphi_n=\varphi,
\]
où $Q_+(h)(x,u)=\E[h((x+uZ)_+)]$.
Si $\varphi$ et $x\mapsto x\sigma(x)$ sont convexes Lipschitz, la convexité se propage aux $\varphi_k$.
Enfin, sous des hypothèses de domination $0\le \varphi\le \psi$ et $\sigma\le \vartheta$, on obtient des inégalités de prix par passage à la limite $n\to\infty$.

\part{Sujet 09 (10 janvier 2013)}

\section{Problème 1 -- Autour du schéma de Milstein}
\subsection{Cadre}
On considère l'équation différentielle stochastique (EDS)
\begin{equation}\label{eq:eds-general}
  dX_t=b(t,X_t)\,dt+\sigma(t,X_t)\,dW_t,\qquad X_0=x,\quad t\in[0,T],
\end{equation}
où $b$ et $\sigma$ sont continues sur $[0,T]\times\R$, à croissance linéaire en $x$ uniformément en $t$.

\subsection{1.a -- Borne de moment quadratique du supremum}
\begin{proposition}\label{prop:sup2}
Sous les hypothèses précédentes (et existence d'une solution forte), il existe $C_{b,\sigma,T}\in(0,\infty)$ tel que
\[
  \E\Big[\sup_{t\in[0,T]}|X_t|^2\Big]\le C_{b,\sigma,T}(1+|x|^2).
\]
\end{proposition}
\begin{proof}
On écrit l'intégrale stochastique
\[
X_t=x+\int_0^t b(s,X_s)\,ds+\int_0^t \sigma(s,X_s)\,dW_s.
\]
Par $(a+b+c)^2\le 3(a^2+b^2+c^2)$,
\[
\sup_{u\le t}|X_u|^2 \le 3|x|^2 + 3\sup_{u\le t}\Big|\int_0^u b(s,X_s)\,ds\Big|^2
 +3\sup_{u\le t}\Big|\int_0^u \sigma(s,X_s)\,dW_s\Big|^2.
\]
Pour le terme de drift, par Cauchy--Schwarz,
\[
\sup_{u\le t}\Big|\int_0^u b(s,X_s)\,ds\Big|^2
\le \Big(\int_0^t |b(s,X_s)|\,ds\Big)^2
\le t\int_0^t |b(s,X_s)|^2ds.
\]
Pour le terme martingale, l'inégalité de Burkholder--Davis--Gundy (BDG) donne l'existence d'une constante universelle $C>0$ telle que
\[
\E\Big[\sup_{u\le t}\Big|\int_0^u \sigma(s,X_s)\,dW_s\Big|^2\Big]
\le C\,\E\Big[\int_0^t |\sigma(s,X_s)|^2\,ds\Big].
\]
En combinant, on obtient
\[
\E\sup_{u\le t}|X_u|^2
\le 3|x|^2 + C\int_0^t \E\big[|b(s,X_s)|^2+|\sigma(s,X_s)|^2\big]\,ds.
\]
Par croissance linéaire, il existe $A>0$ tel que $|b(t,y)|+|\sigma(t,y)|\le A(1+|y|)$, donc
$|b|^2+|\sigma|^2\le C(1+|y|^2)$. D'où
\[
\E\sup_{u\le t}|X_u|^2
\le C(1+|x|^2)+C\int_0^t \E\sup_{r\le s}|X_r|^2\,ds.
\]
La conclusion résulte du lemme de Gronwall.
\end{proof}

\subsection{1.b -- Conditions suffisantes d'existence et d'unicité forte}
\begin{proposition}
Si $b(t,\cdot)$ et $\sigma(t,\cdot)$ sont globalement Lipschitz en $x$ (uniformément en $t\in[0,T]$) et à croissance linéaire, alors \eqref{eq:eds-general} admet une \emph{solution forte unique} sur $[0,T]$.
\end{proposition}
\begin{remark}
Des variantes existent avec des hypothèses locales + non explosion, ou monotonicité; ici on retient le critère standard de cours.
\end{remark}

\subsection{1.c -- Schéma d'Euler à pas $\Delta=T/n$}
On définit le schéma d'Euler discret $(\bar X_{t_i})_{i=0,\dots,n}$ par
\[
\bar X_{t_0}=x,\qquad
\bar X_{t_{i+1}}=\bar X_{t_i}+b(t_i,\bar X_{t_i})\Delta+\sigma(t_i,\bar X_{t_i})\Delta W_i,\quad i=0,\dots,n-1.
\]

\subsection{Rappel -- Schéma de Milstein (scalaire)}
Sous régularité en $x$ de $\sigma$, le schéma de Milstein discret s'écrit
\begin{equation}\label{eq:milstein}
X^{\mathrm{mil}}_{t_{i+1}}
=X^{\mathrm{mil}}_{t_i}+b(t_i,X^{\mathrm{mil}}_{t_i})\Delta+\sigma(t_i,X^{\mathrm{mil}}_{t_i})\Delta W_i
+\frac12\sigma(t_i,X^{\mathrm{mil}}_{t_i})\partial_x\sigma(t_i,X^{\mathrm{mil}}_{t_i})\big((\Delta W_i)^2-\Delta\big).
\end{equation}

\subsection{2.a -- Positivité du Milstein discret (trinôme)}
On suppose désormais:
\begin{itemize}[leftmargin=*]
  \item $b(t,\cdot),\sigma(t,\cdot)$ dérivables en $x$;
  \item pour tout $t$, la fonction $x\mapsto \sigma^2(t,x)$ est croissante sur $\R_+$;
  \item $\sigma$ ne s'annule pas sur $[0,T]\times\R_+$ (donc de signe constant sur ce domaine).
\end{itemize}
On suppose en outre que pour tout $t\in[0,T]$ et $\xi\in\R_+$,
\begin{equation}\label{eq:cond-pos-milstein}
b(t,\xi)-\frac14(\sigma^2)'(t,\xi)\ge 0
\qquad\text{et}\qquad
\frac{\sigma(t,\xi)}{\partial_x\sigma(t,\xi)}\le 2\xi.
\end{equation}
\begin{proposition}
Si $x\ge 0$ et \eqref{eq:cond-pos-milstein} est satisfaite, alors pour tout $i\in\{0,\dots,n\}$,
$X^{\mathrm{mil}}_{t_i}\ge 0$ p.s.
\end{proposition}
\begin{proof}
On procède par récurrence. Supposons $X_i:=X^{\mathrm{mil}}_{t_i}\ge 0$. Écrivons \eqref{eq:milstein} à l'instant $t_i$ en posant
$\sigma_i=\sigma(t_i,X_i)$, $\sigma_i'=\partial_x\sigma(t_i,X_i)$, $b_i=b(t_i,X_i)$, et $u=\Delta W_i$.
On obtient
\[
X_{i+1}=X_i+b_i\Delta+\sigma_i u+\frac12\sigma_i\sigma_i'(u^2-\Delta)
=Au^2+Bu+C,
\]
avec $A=\frac12\sigma_i\sigma_i'$, $B=\sigma_i$ et $C=X_i+b_i\Delta-\frac12\sigma_i\sigma_i'\Delta$.

Comme $x\mapsto \sigma^2(t_i,x)$ est croissante sur $\R_+$, on a $(\sigma^2)'(t_i,X_i)=2\sigma_i\sigma_i'\ge 0$, donc $A\ge 0$.
Si $A=0$ alors $\sigma_i\sigma_i'=0$; avec $\sigma_i\neq 0$ sur $\R_+$, cela impose $\sigma_i'=0$ et \eqref{eq:cond-pos-milstein} donne $\sigma_i/\sigma_i'$ non défini: ce cas est exclu (ou traité séparément) en pratique. On suppose donc $A>0$.

Le minimum du trinôme $Au^2+Bu+C$ sur $\R$ vaut $C-\frac{B^2}{4A}$. Ici
\[
C-\frac{B^2}{4A}=X_i+\Delta\Big(b_i-\frac12\sigma_i\sigma_i'\Big)-\frac{\sigma_i}{2\sigma_i'}.
\]
Or $b_i-\frac12\sigma_i\sigma_i'=b_i-\frac14(\sigma^2)'(t_i,X_i)\ge 0$ par \eqref{eq:cond-pos-milstein}, donc
\[
C-\frac{B^2}{4A}\ge X_i-\frac{\sigma_i}{2\sigma_i'}.
\]
Enfin $\frac{\sigma_i}{\sigma_i'}\le 2X_i$ implique $X_i-\frac{\sigma_i}{2\sigma_i'}\ge 0$. Ainsi $X_{i+1}\ge 0$ p.s.
\end{proof}

\subsection{2.b -- Positivité du Milstein continu}
Le schéma de Milstein \emph{continu} sur $[t_i,t_{i+1}]$ s'écrit (interpolation ``authentique''):
\[
\bar X^{\mathrm{mil}}_t
=X_i+b_i(t-t_i)+\sigma_i(W_t-W_{t_i})+\frac12\sigma_i\sigma_i'\Big((W_t-W_{t_i})^2-(t-t_i)\Big).
\]
Pour $s=t-t_i$ fixé, c'est un trinôme en $u=W_t-W_{t_i}$ dont le minimum vaut
\[
X_i+s\Big(b_i-\frac12\sigma_i\sigma_i'\Big)-\frac{\sigma_i}{2\sigma_i'}\ge 0
\]
par les mêmes arguments que ci-dessus. Donc $\bar X^{\mathrm{mil}}_t\ge 0$ p.s. pour tout $t$.

\subsection{3.a -- Transformation $Y_t=e^{\kappa t}X_t$}
Soit $\kappa\in\R$ et $Y_t=e^{\kappa t}X_t$. Par Itô et la règle de dérivation,
\[
dY_t=\kappa e^{\kappa t}X_t\,dt+e^{\kappa t}dX_t
=\Big(\kappa Y_t+e^{\kappa t}b(t,e^{-\kappa t}Y_t)\Big)dt+e^{\kappa t}\sigma(t,e^{-\kappa t}Y_t)\,dW_t.
\]
Ainsi $Y$ résout une EDS de la forme $dY_t=\tilde b(t,Y_t)dt+\tilde\sigma(t,Y_t)dW_t$ avec
\[
\tilde b(t,y)=\kappa y+e^{\kappa t}b(t,e^{-\kappa t}y),
\qquad
\tilde\sigma(t,y)=e^{\kappa t}\sigma(t,e^{-\kappa t}y).
\]

\subsection{3.b -- Conditions de positivité avec le paramètre $\kappa$}
On souhaite appliquer \eqref{eq:cond-pos-milstein} à $(\tilde b,\tilde\sigma)$.
On calcule
\[
\partial_y\tilde\sigma(t,y)=\partial_x\sigma(t,e^{-\kappa t}y),
\qquad
(\tilde\sigma^2)'_y(t,y)=e^{\kappa t}(\sigma^2)'_x(t,e^{-\kappa t}y).
\]
En posant $x=e^{-\kappa t}y$ (donc $y=e^{\kappa t}x$), les conditions deviennent:
\[
\frac{\tilde\sigma(t,y)}{\partial_y\tilde\sigma(t,y)}\le 2y
\iff
\frac{\sigma(t,x)}{\partial_x\sigma(t,x)}\le 2x,
\]
et
\[
\tilde b(t,y)-\frac14(\tilde\sigma^2)'_y(t,y)\ge 0
\iff
b(t,x)-\frac14(\sigma^2)'(t,x)+\kappa x\ge 0
\quad\text{pour tout }x\ge 0.
\]
Ainsi, un choix de $\kappa$ suffisamment grand (si besoin) peut rétablir la condition de drift.

\subsection{4.a -- Conditions suffisantes de convergence $L^p$ des schémas de Milstein}
Un jeu d'hypothèses standard (suffisant) pour la convergence forte $L^p$ du Milstein sur $[0,T]$ est:
\begin{itemize}[leftmargin=*]
  \item $b$ et $\sigma$ globalement Lipschitz en $x$ (uniformément en $t$), croissance linéaire;
  \item $\sigma$ de classe $C^1$ en $x$, et la fonction $x\mapsto \sigma(t,x)\partial_x\sigma(t,x)$ Lipschitz en $x$ (uniformément en $t$).
\end{itemize}
Alors, pour tout $p\in[1,\infty)$,
\[
\Big\|\sup_{t\le T}|Y_t-\bar Y^{\mathrm{mil}}_t|\Big\|_{L^p}\xrightarrow[n\to\infty]{}0
\quad\text{et typiquement l'ordre fort est }1.
\]

\subsection{4.b -- Positivité de la solution de l'EDS}
Si, pour tout $n$, le schéma $\bar X^{\mathrm{mil},n}$ (ou $\bar Y^{\mathrm{mil},n}$) est positif p.s. sur $[0,T]$ et converge vers la solution $X$ (ou $Y$) dans $C([0,T])$ en probabilité (ou en loi), alors la limite est également positive p.s. puisque l'ensemble
\[
\{ \omega\in C([0,T]) : \forall t,\ \omega(t)\ge 0\}
\]
est fermé pour la norme sup.

\section{Modèles CIR et OU (questions 5--6)}
\subsection{5.a -- Quand les théorèmes standards ne s'appliquent pas}
Pour
\[
dX_t=(a-bX_t)\,dt+\vartheta |X_t|^\alpha\,dW_t,\qquad X_0=x\ge 0,\quad \alpha\in\Big[\frac12,1\Big],
\]
les théorèmes les plus standards d'existence/unicité forte supposent que le coefficient de diffusion $x\mapsto \vartheta|x|^\alpha$ est globalement Lipschitz. Or ce n'est pas le cas au voisinage de $0$ dès que $\alpha<1$. On ne peut donc pas conclure directement par ces résultats.

\subsection{5.b -- Discussion de la positivité via le critère Milstein}
Sur $\R_+$, $\sigma(x)=\vartheta x^\alpha$ et $\sigma'(x)=\vartheta\alpha x^{\alpha-1}$, donc
\[
\frac{\sigma(x)}{\sigma'(x)}=\frac{x}{\alpha}\le 2x \iff \alpha\ge \frac12.
\]
De plus $(\sigma^2)'(x)=2\alpha\vartheta^2 x^{2\alpha-1}$ et la condition de drift (avec paramètre $\kappa$) devient
\[
a-bx+\kappa x-\frac14(\sigma^2)'(x)
a+(\kappa-b)x-\frac{\alpha\vartheta^2}{2}x^{2\alpha-1}\ge 0,\qquad \forall x\ge 0.
\]
Pour $\alpha=\frac12$, le terme $x^{2\alpha-1}$ est constant et l'inégalité à l'origine impose notamment $a\ge \vartheta^2/4$ (condition suffisante ici).

\subsection{6.a -- OU et CIR via la transformation $X_t=Z_t^2$}
Considérons l'OU
\[
dZ_t=-\lambda Z_t\,dt+\sigma\,dW_t,\qquad Z_0=z.
\]
Par Itô,
\[
d(Z_t^2)=2Z_t\,dZ_t+d\langle Z\rangle_t
+(\sigma^2-2\lambda Z_t^2)\,dt+2\sigma Z_t\,dW_t.
\]
En posant $X_t=Z_t^2$ et en notant que $Z_t=\pm\sqrt{X_t}$, on reconnaît une dynamique de type CIR
\[
dX_t=(a-bX_t)\,dt+\vartheta\sqrt{X_t}\,d\widetilde W_t,
\]
avec $a=\sigma^2$, $b=2\lambda$, $\vartheta=2\sigma$ (donc $a=\vartheta^2/4$).

\subsection{6.b -- Euler pour OU}
Avec $\Delta=T/n$,
\[
\bar Z_{t_{i+1}}=(1-\lambda\Delta)\bar Z_{t_i}+\sigma\Delta W_i.
\]

\subsection{6.c -- Bornes uniformes de moments en $L^2$}
Comme $Z_T$ est gaussien (formule explicite), $\E|Z_T|^2<\infty$. De plus, $\bar Z_{t_i}$ vérifie une récurrence affine gaussienne; un calcul de variance récursif fournit
\[
\sup_{n\ge 1}\E|\bar Z_T|^2<\infty,
\]
donc $\sup_n \E(|Z_T|+|\bar Z_T|)^2<\infty$.

\subsection{6.d -- Schéma pour CIR quand $a=\vartheta^2/4$}
Dans le cas $a=\vartheta^2/4$, on peut simuler (ou approximer) $Z_T$ puis poser
\[
\bar X_T=(\bar Z_T)^2,
\]
ce qui donne un schéma pour $X_T$.
Sous les hypothèses de régularité de cours, une erreur faible d'ordre $1/n$ le long des fonctions Lipschitz $f:\R_+\to\R$ se déduit de la régularité du modèle OU et du transfert $X=Z^2$.

\subsection{6.e -- Simulation exacte du CIR quand $a=\vartheta^2/4$}
Simuler exactement $Z_T$ (gaussien explicite) puis retourner $X_T=Z_T^2$.
(Plus généralement, la loi exacte de $X_T$ est une $\chi^2$ non centrale.)

\section{Heston (questions 7--8)}
\subsection{7.a -- Interprétation de $\rho$}
Dans Heston,
\[
dS_t=rS_t\,dt+\sqrt{v_t}\,S_t(\rho\,dW_t+\sqrt{1-\rho^2}\,dB_t),\qquad
dv_t=(a-bv_t)\,dt+\sigma\sqrt{v_t}\,dW_t,
\]
le paramètre $\rho\in[-1,1]$ représente la \emph{corrélation instantanée} entre le bruit $W$ pilotant la variance $v$ et la composante corrélée du bruit du sous-jacent $S$.

\subsection{7.b -- Conditionnement et représentation type Black--Scholes}
On écrit la dynamique de $\log S_t$:
\[
d(\log S_t)=\Big(r-\frac12 v_t\Big)dt+\sqrt{v_t}\big(\rho\,dW_t+\sqrt{1-\rho^2}\,dB_t\big).
\]
Donc
\[
\log S_T=\log S_0+rT-\frac12\int_0^T v_tdt+\rho\int_0^T\sqrt{v_t}\,dW_t
 +\sqrt{1-\rho^2}\int_0^T\sqrt{v_t}\,dB_t.
\]
Conditionnellement à $\cF_T^W=\sigma(W_s,0\le s\le T)$, le terme
$\int_0^T\sqrt{v_t}\,dB_t$ est gaussien centré de variance $\int_0^T v_tdt$, donc $S_T$ est un log-normal conditionnel.
En posant
\[
\Xi=S_0\exp\Big(\rho\int_0^T\sqrt{v_t}\,dW_t-\frac{\rho^2}{2}\int_0^T v_tdt\Big),
\qquad
s=\sqrt{\frac{1-\rho^2}{T}\int_0^T v_tdt},
\]
on obtient, pour un payoff $\varphi$ de croissance polynomiale,
\[
\E\big[e^{-rT}\varphi(S_T)\mid\cF_T^W\big]=\mathrm{PBS}(\Xi,r,s).
\]

\subsection{7.c -- Estimateur de Monte Carlo conditionnel}
Algorithme: simuler $v$ (et donc $W$) puis calculer $\Xi$ et $\int_0^T v_tdt$; remplacer la simulation de $B$ par la formule fermée $\mathrm{PBS}$. C'est un \emph{Rao--Blackwellisation}, donc variance en général plus faible qu'un estimateur ``naïf'' simulant aussi $B$.
La méthode est réellement implémentable si $\varphi$ est un payoff dont le prix Black--Scholes $\mathrm{PBS}$ se calcule (au moins numériquement) de manière stable.

\subsection{8 -- Schéma de Milstein et positivité de $v$}
On met en œuvre un schéma (Milstein ou autre) sur $v$ garantissant $v\ge 0$ (par exemple un Milstein ``positif'' comme étudié plus haut, ou un schéma exact/implicitement biaisé adapté au CIR). Les conditions portent sur le couple $(a,b,\sigma)$ du CIR pour que le schéma reste dans $\R_+$.

\section{Problème 2 -- Autour du pont de diffusion}
\subsection{1.a -- Loi de l'infimum d'un pont brownien}
Pour le pont brownien $Y^{B,T,y}$ allant de $0$ à $y$ en temps $T$, on rappelle la loi du supremum:
\[
\Pp\Big(\sup_{t\le T}Y^{B,T,y}_t\le z\Big)=
\begin{cases}
1-\exp\!\big(-\frac{2}{T}z(z-y)\big), & z\ge \max(y,0),\\
0, & z\le \max(y,0).
\end{cases}
\]
Par symétrie, $-Y^{B,T,y}$ est un pont allant de $0$ à $-y$, donc
\[
\Pp\Big(\inf_{t\le T}Y^{B,T,y}_t\le z\Big)
=\Pp\Big(\sup_{t\le T}(-Y^{B,T,y}_t)\ge -z\Big)
=
\begin{cases}
\exp\!\big(-\frac{2}{T}z(z-y)\big), & z\le \min(y,0),\\
1, & z\ge \min(y,0).
\end{cases}
\]

\subsection{1.b -- Loi de l'infimum d'un pont affine (cas demandé)}
Fixons $x,y\in\R$ et $\lambda>0$. Considérons
\[
X_t=x+\frac{t}{T}(y-x)+\lambda Y_t^{B,T,0}.
\]
C'est un pont gaussien allant de $x$ à $y$ dont la fluctuation est de variance $\lambda^2$.
Pour $z\le \min(x,y)$ on a
\[
\Pp\Big(\inf_{t\le T}X_t\le z\Big)=\exp\!\Big(-\frac{2(x-z)(y-z)}{\lambda^2 T}\Big),
\]
et pour $z\ge \min(x,y)$ la probabilité vaut $1$.

\subsection{1.c -- Simulation par inversion}
Si $U\sim\mathrm{Unif}(0,1)$, posez $d=-(\lambda^2T/2)\log U\ge 0$ puis définissez
\[
Z=\frac{x+y-\sqrt{(x-y)^2+4d}}{2}.
\]
Alors $Z$ a la loi de $\inf_{t\le T}X_t$.

\subsection{2.a -- Loi conditionnelle du schéma d'Euler continu}
Le schéma d'Euler continu associé à
\[
dX_t=b(X_t)\,dt+\sigma(X_t)\,dW_t
\]
sur $[t_i,t_{i+1}]$ s'écrit
\[
\bar X_t=\bar X_{t_i}+b(\bar X_{t_i})(t-t_i)+\sigma(\bar X_{t_i})(W_t-W_{t_i}).
\]
Conditionnellement au squelette $\sigma(\bar X_{t_k},k=0,\dots,n)$, sur chaque intervalle $[t_i,t_{i+1}]$, le processus $(W_t-W_{t_i})_{t\in[t_i,t_{i+1}]}$ conditionné à sa valeur finale est un \emph{pont brownien}. Ainsi, conditionnellement au squelette, $(\bar X_t)$ est un \emph{pont gaussien affine} entre les extrémités $\bar X_{t_i}$ et $\bar X_{t_{i+1}}$.

\subsection{2.b -- MC pour $\E[f(X_T,\inf_{t\le T}X_t)]$ (méthode pont d'Euler)}
Méthode:
\begin{itemize}[leftmargin=*]
  \item Simuler le squelette Euler $(\bar X_{t_i})_{i=0..n}$.
  \item Sur chaque intervalle $[t_i,t_{i+1}]$, conditionnellement à $(\bar X_{t_i},\bar X_{t_{i+1}})$, simuler l'infimum du pont gaussien affine via la formule de la question 1.c (après réduction au cas $x\to y$ et variance $\lambda=\sigma(\bar X_{t_i})$).
  \item Prendre le minimum global des infimums sur les $n$ intervalles.
  \item Moyenner $f(\bar X_T,\inf \bar X)$.
\end{itemize}

\subsection{2.c -- Vitesse de convergence}
Sous hypothèses Lipschitz et si $f$ est Lipschitz, l'erreur due à la discrétisation en temps est contrôlée par la convergence forte du schéma d'Euler: typiquement d'ordre $1/2$, i.e. $O(n^{-1/2})$.

\subsection{3 -- Option barrière et facteur correctif}
Pour une barrière basse $L<x_0$, on veut approximer
\[
\E[g(X_T)\1_{\{\inf_{t\le T}X_t\ge L\}}].
\]
\paragraph{Idée.}
Remplacer $X$ par l'Euler continu $\bar X$ et conditionner par le squelette. Sur chaque intervalle, la probabilité conditionnelle de ne pas franchir $L$ est explicite via le pont brownien, ce qui donne un \emph{poids correctif}.
\medskip

Si $x_{i-1}=\bar X_{t_{i-1}}$ et $x_i=\bar X_{t_i}$ sont $\ge L$, alors
\[
\Pp\Big(\inf_{t\in[t_{i-1},t_i]}\bar X_t\ge L\ \big|\ \bar X_{t_{i-1}}=x_{i-1},\bar X_{t_i}=x_i\Big)
=1-\exp\!\Big(-\frac{2(x_{i-1}-L)(x_i-L)}{\Delta\,\sigma^2(x_{i-1})}\Big).
\]
Ainsi,
\[
\E\big[g(\bar X_T)\1_{\{\inf_{t\le T}\bar X_t\ge L\}}\big]
=\E\Big[g(\bar X_T)\1_{\{\min_{0\le k\le n}\bar X_{t_k}\ge L\}}\prod_{i=1}^n\Big(1-e^{-\frac{2(\bar X_{t_{i-1}}-L)(\bar X_{t_i}-L)}{\Delta\,\sigma^2(\bar X_{t_{i-1}})}}\Big)\Big].
\]
Le produit est un \emph{correctif} car il réintroduit (en moyenne conditionnelle) les franchissements \emph{entre} dates de discrétisation.

\part{Sujet 10 (5 janvier 2015)}

\section{Quasi--Monte Carlo I}
\subsection{1.a -- Box--Müller en quasi--Monte Carlo}
La transformation de Box--Müller: si $U_1,U_2\sim\mathrm{Unif}(0,1)$ indépendantes, alors
\[
Z_1=\sqrt{-2\log U_1}\cos(2\pi U_2),\qquad
Z_2=\sqrt{-2\log U_1}\sin(2\pi U_2)
\]
fournit $(Z_1,Z_2)$ i.i.d. $\mathcal{N}(0,1)$.
\medskip

En quasi--Monte Carlo (QMC), on remplace $(U_1,U_2)$ par une suite à faible discrépance $(u_k)_{k\ge 1}\subset[0,1]^2$ (par ex. Sobol' ou Halton) et on pose
$(U_1,U_2)=(u_k^{(1)},u_k^{(2)})$, ce qui produit une suite quasi--déterministe $(Z_{1,k},Z_{2,k})$.

\subsection{1.b -- Simulation de $\mathcal{N}(0,I_d)$}
Pour $d\ge 1$, on simule $d$ normales indépendantes par transformation composante par composante:
\begin{itemize}[leftmargin=*]
  \item soit via $\Phi^{-1}$: $Z_j=\Phi^{-1}(U_j)$, où $(U_1,\dots,U_d)$ provient d'une suite QMC sur $[0,1]^d$;
  \item soit par blocs de 2 via Box--Müller en dimension 2 répétée.
\end{itemize}

\section{Quasi--Monte Carlo II (randomisation)}
\subsection{2.a -- Définitions: discrépance, variation finie (Hardy--Krause), Koksma--Hlawka}
\begin{definition}[Discrépance étoile]
Pour des points $\xi_1,\dots,\xi_n\in[0,1]^d$, la discrépance étoile est
\[
D_n^*(\xi_1,\dots,\xi_n)
=\sup_{x\in[0,1]^d}\Big|\frac{1}{n}\sum_{k=1}^n \1_{[0,x)}(\xi_k)-\lambda_d([0,x))\Big|,
\]
où $\lambda_d$ est la mesure de Lebesgue.
\end{definition}

\begin{definition}[Variation au sens de Hardy--Krause]
Une fonction $f:[0,1]^d\to\R$ est dite de variation finie (Hardy--Krause) si la somme des variations (au sens de Vitali) de toutes les restrictions de $f$ aux faces de $[0,1]^d$ est finie; on la note $V(f)$.
\end{definition}

\begin{theorem}[Inégalité de Koksma--Hlawka]
Si $f$ est de variation finie sur $[0,1]^d$, alors pour tout $\xi_1,\dots,\xi_n\in[0,1]^d$,
\[
\Big|\frac1n\sum_{k=1}^n f(\xi_k)-\int_{[0,1]^d}f(u)\,du\Big|
\le V(f)\,D_n^*(\xi_1,\dots,\xi_n).
\]
\end{theorem}

\subsection{2.b -- Conclusion pour des intégrales gaussiennes sur $\R^d$}
On ramène une intégrale gaussienne $\int_{\R^d} g(z)\,\varphi_d(z)\,dz$ à une intégrale sur $[0,1]^d$ par transport (par exemple $z=\Phi^{-1}(u)$ composante par composante). Les bornes a priori de type Koksma--Hlawka s'appliquent alors si $f(u)=g(\Phi^{-1}(u))$ est de variation finie; en pratique, la transformation peut dégrader la variation et les bornes deviennent pessimistes en dimension élevée.

\subsection{Randomisation par \emph{shift} modulo 1}
Pour $\xi\in(0,1)^d$, définissons $f_\xi(u)=f(\{u+\xi\})$.
\begin{proposition}
Si $f$ est Riemann-intégrable, alors $f_\xi$ l'est aussi et
\[
\int_{[0,1]^d} f(\{u+\xi\})\,du=\int_{[0,1]^d}f(u)\,du.
\]
\end{proposition}

\subsection{Estimateur randomisé et variance}
Avec $U_k$ i.i.d. uniformes et une suite déterministe $(\xi_\ell)_{\ell=1..n}\subset(0,1)^d$,
\[
\widehat I_{n,M}(f)=\frac1{Mn}\sum_{k=1}^M\sum_{\ell=1}^n f(\{U_k+\xi_\ell\}).
\]
Alors $\widehat I_{n,M}(f)\to \int f$ p.s. quand $M\to\infty$ et
\[
\Var(\widehat I_{n,M}(f))=\frac{\sigma_n(f)^2}{M},
\qquad
\sigma_n(f)^2=\Var\Big(\frac1n\sum_{\ell=1}^n f(\{U+\xi_\ell\})\Big).
\]

\subsection{Cas périodique isotrope}
Si $f$ est $1$-périodique dans chaque direction, alors $f(\{x+\xi\})=f(x+\xi)$. Sous hypothèse ``$V(f)<\infty$'' et invariance de variation par translation, on obtient
\[
\sigma_n(f)\le D_n^*(\xi_1,\dots,\xi_n)\,V(f),
\]
et pour des suites type Sobol'/Halton, $D_n^*\lesssim (\log n)^d/n$, d'où
\[
\Var(\widehat I_{n,M}(f))\lesssim \frac{(\log n)^{2d}}{n^2M}V(f)^2.
\]

\section{Problème -- Couvrir une option digitale (estimation de $f'(K)$)}
\subsection{Cadre}
Soit $X$ solution de
\[
dX_t=b(X_t)\,dt+\sigma(X_t)\,dW_t,\qquad X_0=x_0,
\]
avec $b,\sigma$ Lipschitz. On définit
\[
f(K)=\Pp(X_T\ge K).
\]
\subsection{Préliminaires}
Schéma d'Euler discret:
\[
\bar X_{t_{i+1}}=\bar X_{t_i}+b(\bar X_{t_i})\Delta+\sigma(\bar X_{t_i})\Delta W_i.
\]
On rappelle la convergence forte (cours): pour tout $p\ge 1$,
\[
\Big\|\sup_{t\le T}|X_t-\bar X_t|\Big\|_{L^p}=O(n^{-1/2}).
\]

\subsection{1 -- Inégalités de comparaison de probabilités}
Pour des v.a. $Y,\bar Y$ et un seuil $K$,
\[
\Pp(Y\ge K)-\Pp(\bar Y\ge K)\le \Pp(\bar Y<K\le Y)+\Pp(Y<K\le \bar Y).
\]
De plus, pour tout $\eta>0$ et $p\ge 1$,
\[
\Pp(\bar Y<K\le Y)\le \eta^{-p}\E\big[|Y-\bar Y|^p\big]+\Pp(K\le Y\le K+\eta),
\]
et de même pour $\Pp(Y<K\le \bar Y)$, ce qui donne
\[
|\Pp(Y\ge K)-\Pp(\bar Y\ge K)|
\le \eta^{-p}\E|Y-\bar Y|^p+\Pp(K-\eta\le Y\le K+\eta).
\]

\subsection{2 -- Application à $f(K)-\bar f_n(K)$}
Supposons que $X_T$ admet une densité $p(\cdot)$ bornée par $\Lambda$ sur un voisinage $[K-\eta_0,K+\eta_0]$. Alors
\[
\Pp(K-\eta\le X_T\le K+\eta)\le 2\Lambda\eta
\quad\text{(pour }\eta\le\eta_0\text{).}
\]
Avec $\E|X_T-\bar X_T|^p\le C n^{-p/2}$, on obtient
\[
|f(K)-\bar f_n(K)|\le C\eta^{-p}n^{-p/2}+2\Lambda\eta.
\]
En optimisant (ou en choisissant $\eta=n^{-(1-\delta)/2}$ et $p$ grand), on obtient une vitesse $n^{-\frac12(1-\delta)}$ (pour tout $\delta\in(0,1/2]$).

\subsection{3 -- Dérivabilité de $f$ et approximation de $f'(K)$}
Si $p$ est Lipschitz sur $[K-\eta_0,K+\eta_0]$, alors $f$ est dérivable en $K$ avec
\[
f'(K)=-p(K),
\]
et, pour $\varepsilon\in(0,\eta_0]$,
\[
\left|f'(K)-\frac{f(K+\varepsilon)-f(K)}{\varepsilon}\right|
\le \frac{\varepsilon}{2}[p]_{\Lip}.
\]

\subsection{4 -- Estimateur MC de $f'(K)$ par différence finie}
Pour $n,M\ge 1$ et $\varepsilon>0$, poser
\[
\widehat\kappa_{n}^{\varepsilon,M}
\;:=\;\frac{1}{M\varepsilon}\sum_{k=1}^M
\Big(\1_{\{\bar X_T^{n,k}\ge K+\varepsilon\}}-\1_{\{\bar X_T^{n,k}\ge K\}}\Big),
\]
où $(\bar X_T^{n,k})_{k\ge 1}$ sont i.i.d. de loi $\bar X_T^n$ (simulation Euler).
Le biais se décompose en ``différence finie'' ($\propto \varepsilon$) et ``discrétisation'' ($\propto n^{-\frac12(1-\delta)}/\varepsilon$). La variance vérifie typiquement
\[
\Var(\widehat\kappa_{n}^{\varepsilon,M})
\lesssim \frac{1}{M\varepsilon^2}\Big(\varepsilon\Lambda+C_\delta n^{-\frac12(1-\delta)}\Big).
\]
Un choix équilibré est $\varepsilon_n\asymp n^{-\frac14(1-\delta)}$, puis $M\asymp n^{\frac34(1-\delta)}$.

\part{Fiche de révision (à connaître)}

\section{Outils généraux EDS (à savoir utiliser)}
\begin{itemize}[leftmargin=*]
  \item \textbf{Itô (scalaire)}: si $dX_t=b_t\,dt+\sigma_t\,dW_t$ et $f\in C^2$,
  \[
  df(X_t)=f'(X_t)\,dX_t+\frac12 f''(X_t)\sigma_t^2\,dt.
  \]
  \item \textbf{Itô (vectoriel)}: $dX_t=b_tdt+\Sigma_t dW_t$, alors $df(X_t)=\nabla f^\top dX_t+\frac12 \mathrm{Tr}(\Sigma_t\Sigma_t^\top \nabla^2 f)\,dt$.
  \item \textbf{BDG} (contrôle du $\sup$): pour $M_t=\int_0^t H_s dW_s$, $\E[\sup_{u\le t}|M_u|^p]\le C_p\E[(\int_0^t |H_s|^2ds)^{p/2}]$.
  \item \textbf{Doob} (martingale): $\| \sup_{u\le t}M_u\|_{L^p}\le \frac{p}{p-1}\|M_t\|_{L^p}$ pour $p>1$.
  \item \textbf{Gronwall} (continu/discret): fermer les inégalités intégrales/récurrences.
  \item \textbf{Girsanov} (changement de drift brownien): si $\theta\in L^2$, alors sous $d\Q=\mathcal E(-\int\theta dW)_T\,d\Pp$, $W_t^\Q=W_t+\int_0^t\theta_sds$ est brownien.
  \item \textbf{Cameron--Martin (gaussien)}: pour $Z\sim\mathcal N(0,I_d)$,
  \[
  \E[h(Z)]=\E\big[h(Z+\theta)e^{-\langle\theta,Z\rangle-\|\theta\|^2/2}\big].
  \]
  \item \textbf{Intégration par parties gaussienne}: si $Z\sim\mathcal N(0,1)$ et $H$ régulière,
  $\E[H(Z)Z]=\E[H'(Z)]$ (utile pour greeks).
  \item \textbf{Borel--Cantelli} (passer de $L^p$ à p.s.): si $\sum_n\Pp(A_n)<\infty$ alors $\Pp(A_n\ i.o.)=0$.
  \item \textbf{Existence/unicité forte (critère standard)}: coefficients Lipschitz + croissance linéaire (ou variantes monotones/locales).
\end{itemize}

\section{Schémas numériques (Euler, Milstein) et ordres}
\begin{itemize}[leftmargin=*]
  \item \textbf{Euler discret}:
  $\bar X_{t_{i+1}}=\bar X_{t_i}+b(t_i,\bar X_{t_i})\Delta+\sigma(t_i,\bar X_{t_i})\Delta W_i$.
  \item \textbf{Euler continu (authentique)} sur $[t_i,t_{i+1}]$:
  $\bar X_t=\bar X_{t_i}+b(t_i,\bar X_{t_i})(t-t_i)+\sigma(t_i,\bar X_{t_i})(W_t-W_{t_i})$.
  \item \textbf{Milstein (scalaire)}:
  \[
  X_{t_{i+1}}^{\mathrm{mil}}=X_i+b_i\Delta+\sigma_i\Delta W_i+\tfrac12\sigma_i\sigma_i'((\Delta W_i)^2-\Delta).
  \]
  \item \textbf{Ordres usuels} (coefficients Lipschitz):
  Euler ordre \emph{fort} $1/2$, ordre \emph{faible} $1$; Milstein (scalaire régulier) ordre fort $1$.
  \item \textbf{Astuce $L^p\Rightarrow$ p.s.}:
  si $\E[\sup_{t\le T}|X_t-\bar X_t^n|^p]\le Cn^{-p/2}$ alors pour tout $\alpha<1/2$, choisir $p$ tel que $p(1/2-\alpha)>1$ et appliquer Borel--Cantelli pour obtenir $n^\alpha\sup_{t\le T}|X_t-\bar X_t^n|\to 0$ p.s.
  \item \textbf{Positivité (CIR, etc.)}: écrire la mise à jour comme un trinôme en $\Delta W$ et forcer le minimum $\ge 0$ via des conditions du type
  $b-\frac14(\sigma^2)'\ge 0$ et $\sigma/\sigma'\le 2x$ sur $\R_+$ (ou ajouter un paramètre $\kappa$ via $e^{\kappa t}X_t$).
\end{itemize}

\section{Sensibilités (greeks) : méthodes et pièges}
\begin{itemize}[leftmargin=*]
  \item \textbf{Différences finies}:
  \[
  f'(x)\approx \frac{f(x+\varepsilon)-f(x-\varepsilon)}{2\varepsilon}
  \quad\text{(biais typiquement }O(\varepsilon^2)\text{)}.
  \]
  \item \textbf{Richardson (accélération)}:
  si $D(\varepsilon)=f'(x)+c\varepsilon^2+O(\varepsilon^3)$, alors
  $-\frac13D(\varepsilon)+\frac43D(\varepsilon/2)=f'(x)+O(\varepsilon^3)$.
  \item \textbf{Common Random Numbers (CRN)}: utiliser le \emph{même} bruit pour $x+\varepsilon$ et $x-\varepsilon$ diminue souvent la variance (covariance négative utile).
  \item \textbf{Pathwise derivative} (quand payoff régulier): dériver $X_T(\theta)$ par rapport au paramètre et passer la dérivée sous l'espérance.
  \item \textbf{Likelihood ratio / score function} (quand payoff non régulier): dériver la densité (ou la loi) au lieu du payoff (utile pour digitales).
  \item \textbf{Piège digital}:
  si $x\mapsto \1_{\{X_T^x\ge K\}}$ n'est que Hölder $1/2$ en $L^2$, alors
  $\Var(\text{différence finie})\sim 1/(\varepsilon M)$ : diminuer $\varepsilon$ \emph{augmente} la variance.
  \item \textbf{Règle} : toujours faire un bilan \emph{biais$^2$ vs variance} pour choisir $(\varepsilon,M)$.
\end{itemize}

\section{Réduction de variance (les ``astuces'' classiques)}
\begin{itemize}[leftmargin=*]
  \item \textbf{Rao--Blackwell / MC conditionnel}: remplacer $H$ par $\E[H\mid\mathcal G]$ réduit la variance.
  Ex.: Heston en conditionnant par $\mathcal F_T^W$ donne un prix Black--Scholes conditionnel.
  \item \textbf{Variables de contrôle}: $H-\beta(C-\E[C])$ avec $\beta=\Cov(H,C)/\Var(C)$.
  \item \textbf{Antithétiques}: coupler $Z$ et $-Z$ (efficace si l'intégrande est monotone).
  \item \textbf{Stratification}: découper l'espace en strates et échantillonner conditionnellement.
  \item \textbf{Importance sampling gaussien} (shift): utiliser Cameron--Martin avec paramètre $\theta$; optimiser $\theta$ en minimisant le second moment $g(\theta)$ (convexe, gradient explicite).
\end{itemize}

\section{Pont brownien et barrières}
\begin{itemize}[leftmargin=*]
  \item \textbf{Pont brownien}: $Y_t=B_t-\frac{t}{T}(B_T-y)$.
  \item \textbf{Non-franchissement} (pont $x\to y$ variance $\lambda^2$):
  \[
  \Pp\Big(\inf_{t\le T} \text{bridge}\ge L\Big)
  =1-\exp\!\Big(-\frac{2(x-L)(y-L)}{\lambda^2T}\Big).
  \]
  \item \textbf{Correctif barrière} (Euler): produit des probabilités conditionnelles de non-franchissement sur chaque pas; réduction de biais et souvent de variance.
\end{itemize}

\section{CIR / OU}
\begin{itemize}[leftmargin=*]
  \item \textbf{OU}: $dZ_t=-\lambda Z_tdt+\sigma dW_t$; loi explicite gaussienne.
  \item \textbf{CIR}: $dX_t=(a-bX_t)dt+\vartheta\sqrt{X_t}dW_t$.
  \item \textbf{Lien}: si $X_t=Z_t^2$ alors $X$ est de type CIR avec $a=\sigma^2$, $b=2\lambda$, $\vartheta=2\sigma$, donc $a=\vartheta^2/4$.
\end{itemize}

\section{Quasi--Monte Carlo}
\begin{itemize}[leftmargin=*]
  \item \textbf{Équidistribution}: la mesure empirique converge vers Lebesgue; critère de Weyl.
  \item \textbf{Discrépance étoile} $D_n^*$: mesure d'écart à l'uniformité.
  \item \textbf{Koksma--Hlawka}: $|\frac1n\sum f(\xi_k)-\int f|\le V(f)D_n^*$.
  \item \textbf{Random shift (Cranley--Patterson)}: $\xi_k\mapsto\{\xi_k+U\}$ rend l'estimateur non biaisé et permet IC via $M$ shifts.
  \item \textbf{Scrambling}: randomisation des suites digitales (Sobol') pour réduire corrélations (utile en grande dimension).
  \item \textbf{Attention aux corrélations arithmétiques}: éviter des constructions naïves de points multivariés à partir de sous-suites trop liées (ex. van der Corput).
\end{itemize}

\section{Digitales: $f(K)=\Pp(X_T\ge K)$}
\begin{itemize}[leftmargin=*]
  \item Si $X_T$ a une densité $p$, alors $f'(K)=-p(K)$.
  \item \textbf{Différence finie}: $f'(K)\approx (f(K+\varepsilon)-f(K))/\varepsilon$.
  \item \textbf{Biais/variance}: équilibrer $\varepsilon$, $n$ (discrétisation), et $M$ (MC).
  \item \textbf{Truc barrière/digitale}: remplacer l'indicatrice par une \emph{espérance conditionnelle} via un pont (brownien/Euler) pour lisser et réduire variance.
\end{itemize}

\end{document}

